\section{Counting equivalence and central-group homogeneity}
\label{sec:base-counting}
\label{base:counting}

We now establish the two facts about atom tuples needed to transfer
supports.  With sufficiently few variable names, counting logic cannot
distinguish the zero-charge and unit-charge structures.  Within either
structure, however, the counting type of a short tuple determines its
orbit under the central group $K=\ker\partial$.  Both statements will be
proved in an auxiliary expansion of the structures from
Section~\ref{str:structures}.  We give the quantitative base argument
from~\cite[Section~4, Propositions~11 and~13]{cptCompanion2026} in full.

Throughout this section, write $V=V_L$, $E=E_L$, and let $G=(V,E)$ denote
the undirected box graph, retaining the fixed edge orientations when
using $\partial$.  Thus an element $k\in K$ adds $k_e$ to the scalar
coordinate of every state on $e$, and acts on configurations by the
induced changes of their incident states.

\subsection{The auxiliary expansion and counting types}

Partition the atoms of $A_b$ into classes with labels
\[
 (e,u),\qquad
   e\in E,\quad u\in(\F^2)^{F(e)},
 \qquad\text{and}\qquad
 (v,\eta),\qquad
   v\in V,\quad \eta\in(\F^2)^{F(v)}.
\]
An edge class consists of the three atoms $(e,u,z)$ with its indicated
$e,u$.  A configuration class consists of the configurations at $v$
whose common vector on each face $f\in F(v)$ is $\eta_f$.  The latter
class has $3^{\deg(v)-1}$ members: its scalar coordinates satisfy just
the one divergence equation.  Every class is therefore nonempty, and
its size is independent of $b$.

Choose a linear order of these labels, using the same order for all
charge assignments at a fixed $L$.  Add the binary relation $\preceq$
which compares atoms by their labels, so that its equivalence classes
are exactly the classes just defined.  Also add three binary relations
$\mathsf D_\delta$, for $\delta\in\F$, given by
\[
 \mathsf D_\delta\bigl((e,u,z),(e',u',z')\bigr)
 \quad\Longleftrightarrow\quad
 e=e'\ \text{ and }\ z'-z=\delta.
\]
They are false on pairs involving a configuration atom.  The vector
patterns $u,u'$ need not be equal.  Denote the expansion by $\bar A_b$.
Its vocabulary is still fixed and binary.  Every element of $K$
preserves it, since central translations fix the labels and preserve
raw scalar differences.  The group $H$ is only required to act on the
original input; these additional relations are used in the proof and
are never supplied to a choiceless program.

For a positive integer $h$, let $C^h$ be first-order logic with equality,
ordinary quantifiers, and exact counting quantifiers
$\exists^{=a}x$, using at most $h$ distinct variable names.  The formula
$\exists^{=a}x\,\varphi$ is true precisely when exactly $a$ elements
satisfy $\varphi$, where $a$ is any nonnegative integer.  Formula length
and counting constants are unrestricted.  In particular, when comparing
structures for a fixed $L$, the formulas may depend on $L$.

For an $r$-tuple $\bar a$ with $r\le h$, its $C^h$-type is the set of
all $C^h$ formulas true at $\bar a$ whose free variables belong to its
$r$ designated positions.  The positions have distinct variable names,
but their assigned atoms may repeat.  Types on different structures are
compared using the same designated positions.

\subsection{A large component after deleting few edges}

Write $\delta_G R$ for the edge boundary of a vertex set $R\subseteq V$;
this is separate from the incidence map $\partial$.  We first show that
a small set of deleted edges cannot split the box into small pieces.
The projection step is the three-dimensional discrete Loomis--Whitney
inequality~\cite{loomisWhitney1949}; its short proof below
also gives the finite-box boundary constant needed here.

\begin{lemma}[Boundary estimate]
\label{base:boundary}
Let $c=1-2^{-1/3}$.  For every $L\ge2$ and every nonempty
$R\subseteq\{0,\ldots,L-1\}^3$ with $|R|\le L^3/2$,
\[
 |\delta_G R|\ge c|R|^{2/3}.
\]
\end{lemma}
\begin{proof}
Let $a_{ij},b_{ik},d_{jk}$ be the indicators of the three projections of
$R$ onto coordinate pairs, and let their sizes be $A,B,D$, respectively.
Every point of $R$ contributes to the product of these indicators.
Twice applying Cauchy--Schwarz gives
\begin{align*}
 |R|
 &\le\sum_{i,j,k}a_{ij}b_{ik}d_{jk}\\
 &\le A^{1/2}
      \left(\sum_{i,j}\left(\sum_k b_{ik}d_{jk}\right)^2\right)^{1/2}
 \le (ABD)^{1/2}.
\end{align*}
For the last step, bound each squared inner sum by
$(\sum_k b_{ik}^2)(\sum_k d_{jk}^2)$ and then sum over $i,j$.
At least one projection consequently has size at least $|R|^{2/3}$.
It indexes that many nonempty coordinate lines meeting $R$.  At most
$|R|/L$ of these lines lie wholly in $R$.  Each of the remaining lines
is a path meeting both $R$ and its complement, and hence contains a
boundary edge.  The lines are disjoint, so these edges are distinct.
Finally,
\[
 |\delta_G R|\ge |R|^{2/3}-\frac{|R|}{L}
 \ge (1-2^{-1/3})|R|^{2/3},
\]
where the last inequality uses $|R|\le L^3/2$.
\end{proof}

Set
\begin{equation}
\label{base:M}
 M=\lfloor L^{7/4}\rfloor.
\end{equation}

\begin{lemma}[Persistent large component]
\label{base:giant}
For all sufficiently large $L$, every graph $G\setminus T$ obtained by
deleting at most $6M$ edges has a unique component with more than
$L^3/2$ vertices.  If $T\subseteq T'$ and $|T'|\le6M$, the component
for $T'$ is contained in the component for $T$.
\end{lemma}
\begin{proof}
Sum Lemma~\ref{base:boundary} over the components $R$ of $G\setminus T$
having at most $L^3/2$ vertices.  Every boundary edge belongs to $T$
and is counted at most twice.  Thus
\begin{equation}
\label{base:small-components}
 \sum_R |R|
 \le c^{-3/2}\sum_R|\delta_G R|^{3/2}
 \le c^{-3/2}(2|T|)^{3/2}
 \le c^{-3/2}(12M)^{3/2}=o(L^3).
\end{equation}
Here $\sum_j a_j^{3/2}\le(\sum_j a_j)^{3/2}$ for nonnegative $a_j$,
and $M^{3/2}=O(L^{21/8})=o(L^3)$.  If every component had at most half
the vertices, the left side would be $L^3$, a contradiction for large
$L$.  Two components cannot both contain more than half the vertices.
After further deletion each component is contained in an old component;
its size forces the new large component to lie in the old large one.
\end{proof}

We call the component in Lemma~\ref{base:giant} the \emph{giant
component}.  The next proof uses it to keep the single unmatched charge
away from all assigned atoms.  Its extension step constructs a bijection
of the entire universes, which is what exact counting requires.

\subsection{Counting equivalence of the two charge assignments}

The argument adapts the defect-moving method of the CFI construction
\cite[Section~6]{cfi1992} and the finite-variable counting-game method
of Immerman and Lander~\cite[Counting Quantifiers]{immermanLander1990}.
We construct whole-universe bijections and prove formula preservation
directly, so no game characterization is an additional premise.

\begin{proposition}[Base counting equivalence]
\label{base:equivalence}
For all sufficiently large $L$, the expansions $\bar A_{b^0}$ and
$\bar A_{b^1}$ satisfy the same $C^M$ sentences, where
$b^0=0$, $b^1=\mathbf1_{v_*}$, and $M$ is as in
\eqref{base:M}.
\end{proposition}
\begin{proof}
An assigned edge atom \emph{touches} its edge.  An assigned configuration
atom at $v$ touches every edge incident with $v$.  The edges touched by
at most $M$ assigned positions form a set $T$ of size at most $6M$.

Call paired assignments in the two structures \emph{compatible} if they
have matching block and vector labels and there exist $t\in\F^E$ and
a vertex $w$ in the giant of $G\setminus T$ such that
\begin{equation}
\label{base:defect}
 \partial t=b^1-b^0-\mathbf1_w,
\end{equation}
and each atom in the first assignment is paired with the atom obtained
by adding $t_e$ on every incident scalar coordinate.  The empty
assignments are compatible: take $t=0$ and $w=v_*$.

Compatible assignments preserve atomic formulas.  On any fixed edge,
common addition of $t_e$ preserves raw scalar differences, and also the
corrected differences defining each $\Zrel_\delta$, since vector
coordinates do not change.  It preserves whether the assigned edge
atom is the state specified by an assigned configuration.  Blocks,
sorts, the preorder, and equality are likewise preserved, including
when positions have repeated values.  An assigned configuration cannot
be at $w$: deleting its whole star isolates its vertex, whereas $w$
lies in a component with more than half the vertices.  Thus the shift
at every assigned configuration satisfies the appropriate change of
divergence in \eqref{base:defect}.

Forgetting an assigned position preserves compatibility.  Its touched
edges may disappear from $T$, and Lemma~\ref{base:giant} places the old
giant inside the new one.  Suppose now that at most $M-1$ positions
remain.  We construct a bijection of universes which retains their
pairings and gives a compatible extension for every possible next atom.

Partition the atoms into whole edge blocks and whole configuration
blocks.  For each such block $J$, let $T_J$ be $T$ together with the
edge of $J$, or the star of its vertex, respectively.  Then
$|T_J|\le6M$.  Choose a vertex $w_J$ in the giant of $G\setminus T_J$.
By Lemma~\ref{base:giant}, both $w$ and $w_J$ lie in the giant of
$G\setminus T$.  Choose a signed path vector $p_J$ supported outside
$T$ with
\[
 \partial p_J=\mathbf1_w-\mathbf1_{w_J};
\]
one may traverse a path from $w_J$ to $w$, or use zero if they agree.
The vector $t_J=t+p_J$ agrees with $t$ on $T$ and satisfies
\eqref{base:defect} with $w_J$ in place of $w$.

On the block $J$, map every atom by its $t_J$-shift.  For an edge block
this is plainly a bijection to the corresponding block.  For a
configuration block at $v$, deleting its star isolates $v$, so
$v\ne w_J$.  Therefore
\[
 (\partial t_J)_v=b^1_v-b^0_v,
\]
which is exactly the condition making translation a bijection between
the two configuration blocks.  All their vector constraints stay the
same.  Taking the union over the disjoint blocks gives a bijection of
the whole universes.  On every retained assignment it agrees with the
old pairing because $p_J$ vanishes on $T$.  If the next atom belongs
to $J$, the extended assignment has witness $t_J,w_J$: its touched
set is $T_J$, and all retained coordinates still have the right images.

Induct on formulas to show that compatible assignments agree on every
$C^M$ formula with their designated free positions.  Atomic formulas
were checked above; Boolean connectives preserve agreement.  At a
quantifier, forget the old value, if any, of its bound variable.  At
most $M-1$ positions remain.  The universe bijection just constructed
pairs extensions with compatible assignments, so induction makes the
body true on equally many extensions.  It preserves ordinary
quantifiers and every exact counting quantifier, including count zero.
Apply this conclusion to the empty assignments.
\end{proof}

\subsection{Extending partial central translations}

For homogeneity, both tuples lie in the same expansion $\bar A_b$.
Their coordinates will prescribe a scalar translation on the touched
edges.  We first identify exactly when it extends to an element of $K$.

\begin{lemma}[Partial-flow extension]
\label{base:partial-flow}
Let $T\subseteq E$ and $\sigma\in\F^T$.  There is a vector
$k\in\ker\partial$ with $k|_T=\sigma$ if and only if every component
$U$ of $G\setminus T$ satisfies
\begin{equation}
\label{base:balance}
 R_U:=\sum_{v\in U}\sum_{\substack{e\in T\\e\ni v}}
          \epsilon_{ve}\sigma_e=0.
\end{equation}
If no extension exists, some nonempty component violating
\eqref{base:balance} has
\[
 |U|\le L^3/2,
 \qquad |U|\le(|T|/c)^{3/2}.
\]
\end{lemma}
\begin{proof}
Extend $\sigma$ by zero to all of $E$.  An extension is equivalent to
a vector on $E\setminus T$ having divergence $-\partial\sigma$.
On a connected graph, an incidence vector always has total sum zero.
Conversely, any specified vertex vector of total sum zero is an
incidence vector: choose a spanning tree, give nontree edges value zero,
and remove leaves successively.  The equation at a leaf determines its
tree-edge value uniquely because its coefficient is $1$ or $-1$.
The total-sum condition makes the last root equation hold.  This also
handles a one-vertex component, where the only possible divergence is
zero.  Applied separately to each component of $G\setminus T$, this
argument gives exactly the conditions \eqref{base:balance}.

The sum of all $R_U$ is zero, since the two incidence signs of each
edge cancel.  If one component fails, at least two fail.  Hence one
failing component has at most $L^3/2$ vertices.  Its boundary is
contained in $T$, so Lemma~\ref{base:boundary} yields
$c|U|^{2/3}\le|\delta_G U|\le|T|$.
\end{proof}

\begin{proposition}[Quantitative $K$-homogeneity]
\label{base:homogeneity}
Let $L\ge2$, $b\in\F^V$, and let $r,M$ be positive integers such that
\begin{equation}
\label{base:homogeneity-bound}
 7r+7(6r/c)^{3/2}+1\le M.
\end{equation}
If two $r$-tuples in $\bar A_b$ have the same $C^M$-type, an element
of $K$ maps the first tuple to the second.  Empty tuples lie in the
same $K$-orbit without any width hypothesis.
\end{proposition}
\begin{proof}
Write the tuples as $\bar a,\bar a'$.  The proof first extracts a
consistent scalar shift on the touched edges, then uses a formula to
exclude every obstruction from Lemma~\ref{base:partial-flow}.

Each preorder class is defined with two variable names.  If $h$ is the
number of atoms in strictly earlier classes, its predicate is
\begin{equation}
\label{base:class-predicate}
 \theta_h(x):=\exists^{=h}y\,
          (y\preceq x\land\neg(x\preceq y)).
\end{equation}
Different classes have different $h$ because every class is nonempty.
Equality of types therefore places $a_i$ and $a_i'$ in the same class
for every $i$.  In particular, their blocks and complete vector
patterns agree.  The two tuples touch the same set $T$ of edges, with
$|T|\le6r$.

An \emph{occurrence} on $e\in T$ is either an edge atom in a tuple
position, or the state on $e$ specified by a configuration in a tuple
position.  In the latter case this state is the unique atom $y$ such
that $\Inc(a_i,y)$ holds and $y$ has its specified edge-class label.
The label includes $e$ and the vector pattern, so it is expressed by
one of the predicates \eqref{base:class-predicate}.  Corresponding
occurrences for $\bar a$ and $\bar a'$ consequently have the same
edge and vector coordinates.

For two occurrences on the same edge, their raw scalar difference is
testable using $\mathsf D_\delta$, introducing their uniquely specified
neighbors when necessary.  Such a test uses at most the $r$ tuple
names, two neighbor names, and one further name reused in class
predicates.  Inequality~\eqref{base:homogeneity-bound} supplies these
$r+3$ names.  The raw differences of corresponding occurrences are
therefore equal.  It follows that for every $e\in T$ there is a single
well-defined value $\sigma_e\in\F$: the scalar change from an
occurrence in $\bar a$ to its counterpart in $\bar a'$.  The value is
independent of the chosen occurrence.

Suppose this partial shift $\sigma$ does not extend to $K$.  By
Lemma~\ref{base:partial-flow}, a component $U$ of $G\setminus T$ has
\[
 R_U\ne0,\qquad |U|\le(6r/c)^{3/2}.
\]
We construct a formula distinguishing $\bar a$ from $\bar a'$ within
the stated variable bound.

Choose, for each $v\in U$, an arbitrary configuration atom $g_v^0$ in
$\bar A_b$, and denote its state on an incident edge $e$ by
$h_{v,e}^0$.  Such configurations exist for every $v$.  On each boundary
edge $e\in\delta_G U\subseteq T$, choose one tuple occurrence and
write its state as $o_e^0$.  Let $z(s)$ denote the scalar coordinate of
an edge state $s$; in the following construction it is only used to
specify one of the three relation symbols $\mathsf D_\delta$.

The formula existentially quantifies a configuration variable $g_v$
for each $v\in U$, a state variable $h_{v,e}$ for every edge incident
with $v$, and, when $o_e^0$ comes from a tuple configuration, one
additional neighbor variable $y_e$.  Write $o_e$ for this $y_e$ or for
the selected tuple variable when the occurrence is already an edge
atom.  The conjunction under these quantifiers imposes the following
conditions:
\begin{enumerate}
\item Each $g_v$ and $h_{v,e}$ has the class of its reference atom
      $g_v^0$ or $h_{v,e}^0$, and $\Inc(g_v,h_{v,e})$ holds.
\item On every edge $e$ with both endpoints $v,w\in U$, choose one
      ordering of the endpoints and impose
      \[
       \mathsf D_{z(h_{w,e}^0)-z(h_{v,e}^0)}
                (h_{v,e},h_{w,e}).
      \]
\item On every boundary edge $e$ with endpoint $v\in U$, impose
      \[
       \mathsf D_{z(h_{v,e}^0)-z(o_e^0)}(o_e,h_{v,e}).
      \]
      When $o_e=y_e$, also require its reference edge class and its
      incidence with the chosen tuple configuration.
\end{enumerate}
All class conditions use \eqref{base:class-predicate}.  Thus this is
a finite formula in the expansion's vocabulary, with $\bar a$ as
its free assignment.  Its length and counting constants may depend
on the chosen reference atoms.  The reference configurations and
states witness that it is true at $\bar a$.

The internal-edge condition also applies when $e\in T$ has both
endpoints in $U$.  It does not link that edge to a tuple occurrence;
the two contributions of such an edge to $R_U$ already cancel.  Also,
the reference configurations need not have matching vector patterns
on their common edge.  This is why $\mathsf D_\delta$ was defined
between all states on the same edge, irrespective of their vectors.

If the formula were true at $\bar a'$, compare its state witnesses to
the reference states and put
\[
 d_{v,e}=z(h_{v,e})-z(h_{v,e}^0).
\]
Both configurations at $v$ have divergence $b_v$, so
\[
 \sum_{e\ni v}\epsilon_{ve}d_{v,e}=0.
\]
The constraints on an internal edge give $d_{v,e}=d_{w,e}$.  On a
boundary edge, the tuple occurrence has scalar change $\sigma_e$;
its uniqueness also gives this change when it is represented by
$y_e$.  The boundary constraint therefore forces
$d_{v,e}=\sigma_e$.  Summing the displayed vertex equations over $U$
cancels every internal edge and leaves $R_U=0$, a contradiction.

To count variable names, retain the $r$ tuple names, use $|U|$
configuration names and at most $6|U|$ incident-state names, and use
at most $|\delta_G U|\le6r$ extra occurrence names.  One additional
name can be reused inside every class predicate without capturing
any of these names.  The total is at most
\[
 r+7|U|+6r+1
 \le7r+7(6r/c)^{3/2}+1\le M.
\]
We have obtained a $C^M$ formula distinguishing the tuples, contrary
to their type equality.  Hence $\sigma$ extends to some $k\in K$.
Its action maps each edge occurrence, and thus every tuple
configuration as well, to its counterpart.  For empty tuples, take
$k=0$.
\end{proof}

The support scale from the preceding section is much smaller than the
width available here.  We record the comparison with all parameters
fixed in the order needed later.

\begin{corollary}[Widths at the support scale]
\label{base:width}
Fix positive integers $q,m$ independently of $L$, and let
\[
 N=N_L=\Theta(L^3),\qquad
 s=\left\lceil2L(q\log_3N+1)^2\right\rceil+1,
 \qquad M=\lfloor L^{7/4}\rfloor.
\]
For all sufficiently large $L$, Proposition~\ref{base:equivalence}
holds, Proposition~\ref{base:homogeneity} applies to every tuple length
$1\le r\le(m+4)s$ in either expansion, and
\[
 \max\{4,m+1\}s\le M.
\]
In particular, both expansions are $K$-homogeneous for equal $C^M$
types of tuples of length at most $\max\{2,m\}s$.
\end{corollary}
\begin{proof}
By $N=\Theta(L^3)$,
$s=O_q(L(1+\log L)^2)$.  Uniformly for $1\le r\le(m+4)s$, the left
side of \eqref{base:homogeneity-bound} is
\[
 O_{q,m}\bigl(L^{3/2}(1+\log L)^3\bigr)=o(L^{7/4}).
\]
The floor in $M$ does not affect this eventual inequality.  Also,
$\max\{4,m+1\}s=o(L^{7/4})$.  Increasing $L$ if necessary supplies
Lemma~\ref{base:giant} and Proposition~\ref{base:equivalence} as well.
The threshold may depend on the fixed $q,m$.
\end{proof}
